\section{AI Compute Density: las solicitudes con esperas largas cambian los supuestos de serverless}

Una web request tradicional busca una latencia corta: la CPU ejecuta, la base de datos responde y la solicitud termina pronto. Las aplicaciones de IA suelen llamar a un model provider externo y hacer streaming durante decenas de segundos o incluso varios minutos; en ese lapso, la CPU local puede quedar ociosa. Si cada solicitud sigue ocupando una instancia completa de forma exclusiva, el desperdicio es enorme. El producto de compute que se anunció en clase evolucionó después en Fluid Compute; la documentación oficial actual destaca la reutilización de warm instances, el procesamiento concurrente, la active CPU y un ciclo de vida de ejecución más largo.

\subsection{Las solicitudes CPU-bound e I/O-bound deben planificarse por separado}

Una solicitud \term{CPU-bound} pasa la mayor parte de su tiempo en cómputo local, y el exceso de concurrencia hace que compita por la CPU; una solicitud \term{I/O-bound} pasa la mayor parte de su tiempo esperando a la red, a la base de datos o al modelo, y una concurrencia moderada permite reutilizar la CPU en los intervalos de espera. Si se define la active CPU time de la solicitud $j$ como $a_j$ y su wall time como $w_j$, la proporción de espera es aproximadamente
\[
r_j=1-\frac{a_j}{w_j}.
\]
Una workload con $r_j$ alto es más adecuada para un multiplex seguro; con $r_j$ bajo se necesitan más bien scale-out y aislamiento de recursos.

\lecturefigure{13-compute-density.png}{Las workloads de IA deben elegir entre scale-out y reutilización concurrente según la CPU activity y la I/O wait.}{Subtítulos locales 24:50--28:10; documentación oficial actual de Vercel sobre Fluid Compute.}

\begin{knowledgebox}{Lectura de la figura: compute density no equivale a concurrencia ilimitada}
Una solicitud I/O-bound sigue ocupando memoria, sockets, file descriptors, stream buffers y event-loop capacity. El límite de concurrencia debe determinarse a partir del latency SLO, la memoria, la downstream quota y las noisy-neighbor tests. El planificador necesita observar la active CPU y el wait profile reales, en lugar de escalar solo según el número de solicitudes.
\end{knowledgebox}

\begin{warningbox}{El anuncio en clase y el producto actual no deben confundirse}
La entrevista solo describió la dirección de «aumentar de forma inteligente la compute density según el profile CPU-bound / I/O-bound». La active-CPU billing, la warm instance reuse, la longer execution y el background work que Fluid Compute incorporó después forman parte de la descripción oficial actual del producto, y no deben presentarse retroactivamente como detalles que en la clase de 2025 ya estaban implementados por completo.
\end{warningbox}

\teachervoice{La prioridad de Rauch es el rendimiento: no se puede permitir que, porque un cliente se vuelva popular de repente, un noisy neighbor degrade la calidad de los demás clientes. El objetivo de la compute density no es solo reducir costos, sino disminuir el idle waste manteniendo los estándares de aislamiento y latencia, y trasladar el capacity tuning de las manos del cliente a la plataforma.}

\subsection{Resumen de la sección}

La IA desplaza las web workloads de ráfagas cortas de CPU hacia esperas largas y streaming. El planificador necesita distinguir entre active compute e I/O wait, y reutilizar la warm capacity sin romper el aislamiento.
